% =====================================================================
%  config/environments.tex
%  Environment bersama: pernyataan matematis bergaya amsthm dan kotak OBE.
%  Kelima kotak OBE di bawah ini adalah komponen wajib setiap bab.
% =====================================================================

% --- Environment matematis (Arsitektur_Buku_Ajar.md Bagian 7) -----------
\newtheorem{definition}{Definisi}[chapter]
\newtheorem{theorem}{Teorema}[chapter]
\newtheorem{lemma}{Lemma}[chapter]
\newtheorem{example}{Contoh}[chapter]

% --- Kotak OBE -----------------------------------------------------------
\newtcolorbox{learningoutcome}{
    colback=blue!5!white,
    colframe=blue!75!black,
    title=Tujuan Pembelajaran (Learning Outcomes),
    fonttitle=\bfseries
}

\newtcolorbox{obeactivity}{
    colback=green!5!white,
    colframe=green!75!black,
    title=Aktivitas Rekomendasi,
    fonttitle=\bfseries
}

\newtcolorbox{obereflection}{
    colback=yellow!5!white,
    colframe=yellow!75!black,
    title=Latihan dan Refleksi,
    fonttitle=\bfseries
}

\newtcolorbox{obeassessment}{
    colback=red!5!white,
    colframe=red!75!black,
    title=Asesmen Kompetensi,
    fonttitle=\bfseries
}

\newtcolorbox{competencychecklist}{
    colback=gray!5!white,
    colframe=gray!75!black,
    title=Checklist Kompetensi,
    fonttitle=\bfseries
}
